\scanpage{232}

This implies the

\uline{Corollary 1} Let us suppose that $u$ commutes with
$\widehat\varphi:\widehat\Pi_{0,3}\longrightarrow
\widehat{\mathcal G}_{0,3}^{+}$ up to inner automorphisms. Then it is possible
\uline{in a unique way} to describe \unsure{$u$|$\varphi$} by a triple
$(\alpha,\beta,f)$ (a priori determined modulo a transformation of type (8),
in such a way that we have (putting $\xi=\alpha\rho(\alpha)^{-1}$,
as usual) \unsure{a} relation
\[
\boxed{\widehat\varphi(\xi)=\beta\sigma_\infty^i}
\qquad\text{for }i\in\mathbb Z/2\mathbb Z\ \text{suitable}.
\tag{10}
\]
\margin{We should examine whether this definition coincides with that on
\unsure{the} page\,\unsure{410'}\,,\,\unsure{413|417}, formulated without
involving $\widehat\varphi$.}
Such a choice \unsure{will be} called \uline{normalized}.

Thus the $u\in\mathcal M_{0,3}^{!\sim}$ which commute with
$\widehat\varphi$ up to inner automorphisms correspond one-to-one to the
\deleted{pairs} triples \deleted{$(\alpha,\beta,f)$}
\deleted{satisfying (10)}
\[
(\alpha,\beta,f),\qquad
\alpha\in\unsure{\widetilde{\widehat\Pi}_{0,3}},\quad
\beta,f\in\widehat\Pi_{0,3}
\]
satisfying, on the one hand the loop relation
\[
\xi\unsure{\doteq|=}\eta\ell_1^\nu
\qquad
\begin{gathered}
\text{where }\xi=\alpha\rho(\alpha)^{-1},\quad
\eta=\beta\sigma_\infty(\beta)^{-1},\\
\text{and where }\nu\in\widehat{\mathbb Z}\ \text{suitable}\\
(\text{such that }2\nu+1\in\widehat{\mathbb Z}^{*})
\end{gathered}
\tag{11}
\]
on the other hand relation (10) (for suitable $i\in\mathbb Z/2\mathbb Z$)
--- neither relation moreover involving $f$.
\margin{231}

\scanpage{233}
\mspage{417}

Let us note that $i\in\mathbb Z/2\mathbb Z$ is defined in terms of
$\xi$ (taking into account
$\beta\in\widehat\Pi_{0,3}=\operatorname{Ker}
(\widehat{\mathcal G}_{0,3}^{+}\to\unsure{\mathcal G_3|\mathfrak S_3})$)
by reducing (10) modulo $\widehat\Pi_{0,3}$. Denoting by
\[
\widehat\varphi^{\bullet}:
\widehat\Pi_{0,3}\longrightarrow\unsure{\mathcal G_3|\mathfrak S_3}
\tag{12}
\]
the composite
\[
\widehat\Pi_{0,3}\xrightarrow{\widehat\varphi}
\widehat{\mathcal G}_{0,3}^{+}
\overset{\mathrm{can}}{\longrightarrow}
\unsure{\mathcal G_3|\mathfrak S_3},
\]
formula (10) also gives
\[
\widehat\varphi^{\unsure{\bullet|\circ}}(\xi)=\sigma_\infty^i
\tag{13}
\]
which does determine $i$\inserted{$\in\mathbb Z/2\mathbb Z$} in terms of
$\xi\in\widehat\Pi_{0,3}$, and implies an intrinsic condition on
the element $\xi$ of $\widehat\Pi_{0,3}$, namely that its image in
$\unsure{\mathcal G_3|\mathfrak S_3}$ under
$\widehat\varphi^{\unsure{\bullet|\circ}}$ should indeed be in the subgroup
$\{1,\sigma_\infty\}$ of $\unsure{\mathcal G_3|\mathfrak S_3}$.
In fact, the kernel of $\widehat\varphi^{\unsure{\bullet|\circ}}$ is
the \unsure{inverse} subgroup of $\widehat\Pi_{0,3}$ generated by
$\ell_0^2,\ell_1^2,\ell_\infty^3$, and the condition being considered
therefore means that modulo this subgroup, $\xi$ should be of the form $\ell_1^i$.

By formula (10), we can eliminate $\beta$ from the datum
$(\alpha,\beta,f)$, and we can replace the datum $\alpha$ by the datum
$\xi$. Thus:
\margin{232}

\scanpage{234}

\uline{Corollary 2} \unsure{The} elements of the \unsure{$^{s}/$group} of
$\mathcal M_{0,3}^{!\sim}$ formed by the $u$ which commute with
$\widehat\varphi$ up to inner automorphisms, is \deleted{\illegible{1}} in one-to-one correspondence with
the set of pairs $(\xi,f)$ \deleted{\unsure{$\in$}}
$(\xi,f\in\widehat\Pi_{0,3})$ satisfying the conditions
(\unsure{not|by} involving anything but $\xi$)
\[
\begin{array}{c@{\qquad}c}
\begin{gathered}(14)\\[1.8ex](15)\end{gathered}
&
\boxed{\begin{gathered}
\xi\rho(\xi)\rho^2(\xi)=1\\[1.8ex]
\xi=\bigl(\widehat\varphi(\xi)\,
\sigma_\infty(\widehat\varphi(\xi))^{-1}\bigr)\ell_1^\nu
\end{gathered}}
\end{array}
\]
\[
(\nu\in\widehat{\mathbb Z}\ \text{suitable such that }
2\nu+1\in\widehat{\mathbb Z}^{*}).
\]
\margin{we also need the ``preliminary'' ``condition''\ldots}

For such a $\xi$, $\widehat\varphi^{\unsure{\bullet|\circ}}(\xi)$
in $\unsure{\mathcal G_3|\mathfrak S_3}$ is of the form
$\sigma_\infty^i$, with $i\in\mathbb Z/2\mathbb Z$, and putting
\[
\left\{\begin{aligned}
&\beta=\widehat\varphi(\xi)\sigma_\infty^i
  \quad(\in\widehat\Pi_{0,3})\\
&\alpha\in\widehat\Pi_{0,3}\quad
  \text{the unique element of }\widetilde{\widehat\Pi}_{0,3}
  =\widehat\Pi_{0,3}\cdot\{1,\tau_\infty\}\\
&\hspace{5em}\text{such that }\alpha\rho(\alpha)^{-1}=\xi
\end{aligned}\right.
\tag{16}
\]
\margin{in fact, we do not really need $\alpha,\beta$}
$u$ is obtained from $\xi,f$ by the formulas
\[
\left\{\begin{aligned}
u(\sigma_\infty)
 &=\operatorname{int}(f^{-1}\unsure{\beta})(\sigma_\infty)\\
 &=\operatorname{int}(f^{-1})(\operatorname{int}(\unsure{\beta})(\sigma_\infty))\\
 &=\operatorname{int}(f^{-1})(\xi\ell_1^{-\nu}\sigma_\infty)\\
u(\rho)
 &=\operatorname{int}(f^{-1}\alpha)(\rho)\\
 &=\operatorname{int}(f^{-1})(\operatorname{int}(\alpha)(\rho))\\
 &=\operatorname{int}(f^{-1})(\unsure{\xi}\rho)
\end{aligned}\right.
\tag{17}
\]

We shall note that in these formulas, $f$ enters only through
composition with $\operatorname{int}(f)$. If we
\deleted{put}\inserted{denote}, for $\xi\in\widehat\Pi_{0,3}$
\margin{233}

\scanpage{235}
\mspage{418}

\margin{This \unsure{notation} is already in use for the
general \unsure{considerations} of the preceding \S
$\unsure{u_{\xi,\eta}}$ depends in \unsure{any} case only on $\xi$\ldots}
satisfying (14)(15), by $u_\xi$ the automorphism
\deleted{$u_\xi(\sigma_\infty)$} corresponding, in terms of
Corollary 2, to the pair $(\xi,1)$, so that
\[
\left\{\begin{aligned}
u_\xi(\sigma_\infty)
 &=\deleted{\unsure{\operatorname{int}(\alpha)}}
 \overbrace{(\xi\ell_1^{-\nu})}^{\unsure{\eta|\nu}}\sigma_\infty\\
u_\xi(\rho)&=\xi\rho
\end{aligned}\right.,
\tag{18}
\]
formulas (17) can also be written, simply
\[
u_{\xi,f}=\operatorname{int}(f^{-1})\,u_\xi.
\tag{19}
\]

On the other hand formulas (5), (6) give here
\[
\widehat\varphi\circ u_\xi=u_\xi\circ\widehat\varphi
\tag{20}
\]
where, of course, on the first \unsure{side|side} $u_\xi$ denotes the
restriction of $u_\xi$ to $\widehat\Pi_{0,3}$.

Let us note

\uline{Corollary 3} The loop automorphisms of
$\widehat{\unsure{\mathcal G}}_{0,3}^{+}$ of the form $u_\xi$
\inserted{($\xi$ \unsure{normalized|normalized}$^*$)} are exactly those for
which we have the \deleted{properties} formulas (20) (i.e. which
\deleted{\illegible{1}} commute with $\widehat\varphi$) and which moreover
normalize $L_0$.
\margin{$^*$ in the sense of page\,416'}
\figure{\begin{tikzpicture}[x=1mm,y=1mm,line width=.3pt,line cap=round]
\draw (0,35) .. controls (-.16,29) and (.15,23) .. (-.02,17);
\draw (.65,35.0) -- (.88,34.1) -- (1.38,34.55);
\draw (0,9.5) .. controls (-.08,6.3) and (-.16,2.7) .. (0,0);
\end{tikzpicture}}
\margin{234}

\scanpage{236}

Indeed, in terms of the \inserted{triples} $(\alpha,\beta,f)$ \inserted{``special''} these properties are expressed by the \inserted{\unsure{simple}} conditions \deleted{that we have} $f$ of the form $\ell_0^n$ (which expresses the condition of normalizing $L_0$) and that $\lambda$ given by (6) \unsure{case|for} $r=0$, i.e.
\[
\lambda=\widehat\varphi(f)^{-1}f
\tag{21}
\]
be equal to $1$, i.e. the condition
\[
\widehat\varphi(f)=f,
\tag{22}
\]
which in this case reads
\[
\varepsilon_0^n=\underbrace{\ell_0^n}_{\varepsilon_0^{2n}}
\qquad\text{i.e.}\quad \varepsilon_0^n=1
\]
which implies $\unsure{f|f_0}=\ell_0^n=1$, hence $u$ of the form $u_{(\alpha,\beta,1)}=u_\xi$, where $\xi=\alpha\rho(\alpha)^{-1}$.

\uline{Corollary 4} The $u_\xi$ \inserted{normalized} form a \uline{subgroup} of $\operatorname{Norm}_{\mathcal M_{0,3}^{!\sim}}(L_0)$, i.e. $\mathcal M_{0,3}^{!\approx}(0)$, whose \deleted{\unsure{image}}\inserted{intersection} with $\widehat\Pi_{0,3}$ is reduced to $\{1\}$, and whose image in $\mathcal N_{0,3}^{!\sim}$ is \inserted{the subgroup $\mathcal N_{0,3}^{!\approx}$} formed by the outer loop automorphisms of $\widehat{\mathcal G}_{0,3}^{+}$, stabilizing $\widehat\Pi_{0,3}$ (i.e. $u\in\unsure{\mathcal M_{0,3}^{+\widehat{\ }}}$) such that $\widehat\varphi:\widehat\Pi_{0,3}\to\widehat{\mathcal G}_{0,3}^{+}$ commutes with $u$ up to inner automorphisms. We therefore have an isomorphism
\[
\mathcal M_{0,3}^{!\approx}(\unsure{0})
=\{u_\xi\mid \xi\in\widehat\Pi_{0,3}\text{ solution of }(14)(15)\}
\xrightarrow{\sim}\mathcal N_{0,3}^{!\approx}.
\tag{\unsure{23}}
\]
\margin{235}

\scanpage{237}
\mspage{419}

It only remains to check that a $u_\xi$ can be an inner automorphism only if $u_\xi=1$ i.e. $\xi=1$ (hence $\nu=0$). Now $u_\xi$ is a $u_{\alpha,\beta}=u_{\alpha,\beta,1}$ ($\alpha,\beta$ defined \deleted{by} via $\xi$ by $\alpha\rho(\alpha)^{-1}=\xi$, $\beta\sigma_\infty(\beta)^{-1}=\xi \ell_1^{-\nu}$), which is an inner automorphism iff we have $\alpha=\ell_0^n$, $\beta=\ell_0^n$ (cf. proposition on page 404), i.e. $\xi=\ell_0^n \ell_1^{-n}$, $\eta=\beta\sigma_\infty(\beta)^{-1}=\ell_0^n \ell_1^{-n}$, and \deleted{the relation} \deleted{\unsure{for}} $\xi=\eta \ell_1^\nu$ reads $\ell_1^\nu=1$ i.e. $\nu=0$. Relation (15) therefore becomes (since
\[
\begin{gathered}
\widehat\varphi(\xi)=\widehat\varphi(\ell_0^n \ell_1^{-n})=\varepsilon_0^n\sigma_\infty^n\\
\sigma_\infty(\widehat\varphi(\xi))=\varepsilon_1^n\sigma_\infty^n
\end{gathered}
\]
) \deleted{\unsure{gives}} the relation
\[
\ell_0^n \ell_1^{-n}=\varepsilon_0^n\sigma_\infty^{-n}\varepsilon_1^n\sigma_\infty^n,
\]
i.e. (since $\ell_0=\varepsilon_0^2$, $\ell_1=\varepsilon_1^2$) if $n$ even
\[
\varepsilon_0^n=\deleted{\unsure{\sigma_\infty}}\varepsilon_1^{3n}
\qquad\text{whence }n=0,
\]
and if $n$ odd
\[
\varepsilon_0^n=\deleted{\unsure{\varepsilon_0}}\underbrace{\sigma_\infty\varepsilon_1^n\sigma_\infty}_{\sigma_\infty(\varepsilon_1)^n=\varepsilon_0^n}\varepsilon_1^{2n}
\]
i.e. $\varepsilon_1^{2n}=0$, i.e. again $n=0$, q.e.d.

\margin{\uline{N.B.} \illegible{1} $\unsure{\mathcal M_{0,3}^{!\sim}[0]}$ --- \unsure{a} \illegible{2} $\mathcal N_{0,3}^{!\sim}$ in its entirety, \unsure{but it} remains to be seen whether \unsure{the other}\ldots\ \illegible{1}}

\uline{\unsure{Scholium}} We have thus found a canonical partial section of $\mathcal M_{0,3}^{!\sim}$ over the subgroup $\mathcal N_{0,3}^{!\approx}$ of $\mathcal N_{0,3}^{!\sim}$. Since
\margin{236}